% Appendix H, the annotation guide and the review packet (restructure plan, part 3), corrected to the code
% (src/tracecontract/schemas.py, Candidate.final_label and human_disagreement; src/tracecontract/review.py). Sources:
% docs/annotation_guide.md (v0.4); docs/decision_log.md D43, D44, D46, D64. Every number is a macro (\dv).
\section{Annotation guide and review packet}
\label{app:annotation}
No human label exists: the procedure below has not been used yet. The guide for human reviewers is
\texttt{docs/annotation\_guide.md} (v0.4). It covers:
\begin{itemize}
\item contract review: evidence, kind, status, supersession, observability;
\item candidate labels, made from a blinded review packet (\texttt{tracecontract review-packet}, described below);
\item label values: valid, invalid with the violated requirement ids or \texttt{REGRESSION}, or ambiguous;
\item \dv{annotation.reviewers.w} reviewers per task, each under a pseudonymous id (\texttt{R1}, \texttt{R2}, \dots), who
  each record their contract reviews, their labels and an origin guess per candidate in their own label file for the
  task;
\item merging once every task has \dv{annotation.reviewers.w} complete reviews (\texttt{tracecontract merge-reviews}),
  and adjudicating disagreements: a candidate whose reviewers give different label values, or the same value with
  different violated requirements, has no final label until an adjudicated label is added (D44, D46);
\item the review queue and the evidence that a reviewer may consult only after the labels are merged.
\end{itemize}
Per task, the packet holds the user turns of the session, the reference contract, and each held-out candidate's change
and file state under an unguessable random id, in shuffled order. Candidates with identical diffs share an id. The
packet names no candidate id, origin, description, proposal or verifier outcome, nor a fresh attempt's configuration or
variant, and the key that maps ids to candidates stays with a coordinator (D43). No file of the packet marks which
candidates are fresh attempts, but the guide, which the reviewers read, says that a packet may hold them among a task's
own. The packet hides the origin field but not the origin, which a reviewer can still infer from the content: the user's
corrections describe the agent's first attempt, near misses resemble each other and the reference, and in
\texttt{tinydb} the diff structure differs by origin. The
merge reports the share of origin guesses that are right, as a measure of how blind the labels are; no guess names a
fresh attempt's origin, so for the fresh attempts it counts the guesses by configuration and variant instead. The guesses
are never used as labels. Every analysed fresh attempt is queued for the packet (protocol~\S15.4); for the natural arm no
packet has been prepared (D64, decision~8).
